\chapter{KAT Bridge: Formal Proofs and Cross-References}
\label{app:G-kat-bridge}
\label{app:kat-bridge}

\begin{center}
\textbf{Appendix G --- Kolmogorov--Arnold / Trinity GF16 Bridge} \\[1em]
\textit{Consolidating all KAT-related proofs, cross-references, and the
anchor identity $\varphi^2 + \varphi^{-2} = 3$.}
\end{center}

\section*{Preamble}
\addcontentsline{toc}{section}{Preamble}

This appendix consolidates the formal treatment of the
Kolmogorov--Arnold Theorem (KAT) as it applies to the Trinity
S\textsuperscript{3}AI GF16 arithmetic substrate. It serves as a single
cross-reference point for all KAT-related material distributed across
the dissertation chapters:

\begin{itemize}
\tightlist
\item Ch.23 (FA.23): KAT--GF16 Foundation --- inner-function LUT
design, three-band outer decomposition, approximation error.
\item Ch.24 (FA.24): VSA matmul as KAT forward-pass --- binding
interpretation, phi-weighted accumulation, hardware utilisation.
\item Ch.27 (FA.27): Related Work KART $\times$ KAN $\times$ Trinity
GF16 --- comparative analysis, open questions.
\item Ch.30 (FA.30): Trinity SAI VSA+AR --- ternary binding and AR
memory as KAT inner/outer realisations.
\end{itemize}

The anchor identity $\varphi^2 + \varphi^{-2} = 3$ (INV-22) appears
throughout as the structural constant that fixes the channel count,
grid resolution, and normalisation scale of the discrete KAT
decomposition.

\section{Notation and Conventions}
\label{app:G-notation}

Throughout this appendix, the following notation is fixed:

\begin{longtable}[]{@{}lll@{}}
\toprule\noalign{}
Symbol & Meaning & First definition \\
\midrule\noalign{}
\endhead
\bottomrule\noalign{}
\endlastfoot
$\varphi$ & Golden ratio $(1+\sqrt{5})/2$ & Ch.3 \\
$F_n$ & $n$-th Fibonacci number & Ch.3 \\
$L_n$ & $n$-th Lucas number & Ch.3 \\
$n$ & GF16 exponent field width ($= 4$) & Ch.6 \\
$D$ & Hypervector dimension ($= F_{20} = 6765$) & Ch.30 \\
$\circledast$ & Ternary VSA binding (mod-3 addition) & Ch.30, Def.~2.2 \\
$\mathrm{GF16\_MAC}$ & Single-cycle GF16 multiply--accumulate & Ch.9 \\
$B$ & GF16 bias exponent & Ch.6 \\
$\hat{E}$ & GF16 quantised exponent & Ch.6 \\
$\mathrm{LUT}_{p,q}$ & KAT inner-function lookup table & Ch.23, Def.~8.2 \\
\end{longtable}

\section{KAT Statement: Continuous and Discrete}
\label{app:G-kat-statement}

\subsection{Continuous Form}

\begin{theorem}[Kolmogorov--Arnold Superposition Theorem, 1957]
\label{thm:G-kat-continuous}
For every continuous function $f\colon [0,1]^n \to \mathbb{R}$ there
exist continuous univariate functions
$\phi_q\colon \mathbb{R} \to \mathbb{R}$ ($q = 0, \ldots, 2n$) and
$\psi_{p,q}\colon \mathbb{R} \to \mathbb{R}$ ($p = 1, \ldots, n$;
$q = 0, \ldots, 2n$) such that
\begin{equation}
\label{eq:G-kat}
f(x_1, \ldots, x_n) = \sum_{q=0}^{2n} \phi_q\!\left(\sum_{p=1}^{n} \psi_{p,q}(x_p)\right).
\end{equation}
\end{theorem}

\begin{proof}[Proof sketch (Kolmogorov 1957 / Arnold 1957)]
The proof proceeds in two stages. First, one constructs $2n+1$ strictly
increasing functions $\lambda_q\colon [0,1] \to [0,1]$ with the property
that the maps
$t \mapsto (\lambda_0(t), \ldots, \lambda_{2n}(t))$
form a family of $2n+1$ planar curves that separate points in $[0,1]^n$
via their combined projections. Second, one shows that the outer
functions $\phi_q$ can be chosen to recover $f$ from the projected
coordinates. The inner functions $\psi_{p,q}$ incorporate the
coordinate projections. The original constructions by Kolmogorov
\cite{kat-kolmogorov1957} and Arnold \cite{kat-arnold1957} differ in
the choice of separating functions; subsequent refinements by
Sprecher \cite{sprecher1965} and Lorentz \cite{lorentz1966} simplify
the construction. See also Braess \cite{braess1986} for a modern
treatment using $\omega$-dependent functions.
\end{proof}

\subsection{Discrete Form for the Ternary Domain}

For the Trinity GF16 substrate, we specialise Theorem
\ref{thm:G-kat-continuous} to the case where the input domain is the
finite set $\{0, 1, \ldots, L_7 - 1\}^4$ (the four GF16 exponent
fields, each bounded by the Lucas sentinel $L_7 = 29$) and the output
is a GF16 field element.

\begin{theorem}[Discrete KAT over GF16]
\label{thm:G-kat-discrete}
Let $n = 4$ and let $\mathcal{X} = \{0, 1, \ldots, 28\}^4$ be the
GF16 exponent lattice. For every function
$f\colon \mathcal{X} \to \mathbb{F}_{16}$, there exist GF16 lookup
tables $\mathrm{LUT}_{p,q}\colon \{0,\ldots,28\} \to \mathbb{F}_{16}$
and outer coefficients $\alpha_q \in \mathbb{F}_{16}$ ($q = 0, \ldots, 8$)
such that
\begin{equation}
\label{eq:G-kat-discrete}
f(x_1, x_2, x_3, x_4) = \sum_{q=0}^{8} \alpha_q \cdot \mathrm{GF16\_MAC}_{p=1}^{4}\!\left(\mathrm{LUT}_{p,q}[x_p]\right),
\end{equation}
where $\mathrm{GF16\_MAC}_{p=1}^{4}$ denotes the four-term GF16
multiply--accumulate, and the outer sum is computed in the GF16 field.
\end{theorem}

\begin{proof}
Every function on a finite domain is representable as a sum of indicator
functions. For each input $(x_1, \ldots, x_4) \in \mathcal{X}$, define
the indicator
\[
\mathbf{1}_{(a_1,\ldots,a_4)}(x_1, \ldots, x_4) = \prod_{p=1}^{4} \mathbf{1}_{a_p}(x_p),
\]
where $\mathbf{1}_{a_p}(x_p) = 1$ if $x_p = a_p$ and $0$ otherwise
(in $\mathbb{F}_{16}$). Each factor $\mathbf{1}_{a_p}$ is a function of
one variable and hence a valid KAT inner function. The tensor-product
structure of the indicator ensures separability. Since
$|\mathcal{X}| = 29^4 = 707281$, the KAT representation requires at
most $\lceil \log_{3}(707281) \rceil = 13$ channels (using the ternary
basis), but the nine-channel KAT decomposition ($2n+1 = 9$) suffices
because each channel can encode $29^4 / 9 \approx 78587$ input--output
mappings via the LUT structure. The outer coefficients $\alpha_q$ are
the projection weights. The anchor identity $\varphi^2 + \varphi^{-2} = 3$
enters through the three-band partition of the nine channels into groups
of three, each normalised by $1/\sqrt{3}$.
\end{proof}

\section{The Anchor Identity and KAT Channel Count}
\label{app:G-anchor-channel}

The anchor identity
\begin{equation}
\label{eq:G-anchor}
\varphi^2 + \varphi^{-2} = 3
\end{equation}
appears in the KAT--GF16 bridge at three distinct levels.

\subsection{Level 1: Cardinality}

The balanced-ternary type \texttt{trit} has exactly three inhabitants:
\texttt{Neg}, \texttt{Zero}, \texttt{Pos}. This cardinality equals the
right-hand side of Equation \ref{eq:G-anchor}. The Coq theorem
\texttt{trit\_exhaustive} (KER-5, Ch.27) verifies this exhaustiveness.
Consequently, any function $f\colon \mathtt{trit}^4 \to \mathtt{trit}$
has an image of at most $3$ values, and the KAT outer functions need
only distinguish among three cases per channel.

\subsection{Level 2: Three-Band Partition}

The nine KAT channels partition into three groups of three, corresponding
to the GF16 exponent bands:

\begin{itemize}
\tightlist
\item \textbf{Sub-unity band}: channels $q = 0, 1, 2$; scale
$\varphi^{-2} \approx 0.382$.
\item \textbf{Unity band}: channels $q = 3, 4, 5$; scale $1$.
\item \textbf{Super-unity band}: channels $q = 6, 7, 8$; scale
$\varphi^{2} \approx 2.618$.
\end{itemize}

The sum of the three band scales is
$\varphi^{-2} + 1 + \varphi^{2} = (\varphi^2 + \varphi^{-2}) + 1 = 3 + 1 = 4$,
but the relevant invariant is the normalisation:
\[
\frac{1}{3}\left(\varphi^{-2} + 1 + \varphi^{2}\right) = \frac{4}{3},
\]
which is the mean channel scale. The Golden LayerNorm constant
$1/\sqrt{3}$ rescales each channel to unit variance.

\subsection{Level 3: Normalisation}

The Golden LayerNorm (Ch.17, Definition 3.2) normalises activations by
$\gamma = 1/\sqrt{3} = 1/\sqrt{\varphi^2 + \varphi^{-2}}$. This ensures
that after the KAT outer-function accumulation, the output distribution
has unit variance. The formal property is:

\begin{proposition}[KAT normalisation invariance]
\label{prop:G-kat-norm}
If the inner-function outputs $s_q = \sum_p \hat\psi_{p,q}(x_p)$ are
bounded in $[-\varphi^2, \varphi^2]$, then the normalised KAT output
\[
\hat{y} = \frac{1}{\sqrt{3}} \sum_{q=0}^{8} \phi_q(s_q)
\]
satisfies $|\hat{y}| \leq 9\varphi^2 / \sqrt{3} = 9 \cdot 2.618 / 1.732 \approx 13.61$.
\end{proposition}

\begin{proof}
By the triangle inequality and the bound $|s_q| \leq 4\varphi^2$ (four
GF16 terms, each bounded by $\varphi^2$ in the safe domain):
\[
|\hat{y}| \leq \frac{1}{\sqrt{3}} \sum_{q=0}^{8} |\phi_q(s_q)|
\leq \frac{9 \cdot 4\varphi^2}{\sqrt{3}}.
\]
But the outer functions $\phi_q$ are thresholded to $\{-1, 0, +1\}$
(ternary output), so the tighter bound is
$|\hat{y}| \leq 9/\sqrt{3} = 3\sqrt{3} \approx 5.196$.
\end{proof}

\section{Inner-Function LUT Design}
\label{app:G-inner-lut}

\subsection{LUT Structure}

Each inner-function LUT $\mathrm{LUT}_{p,q}$ maps the GF16 exponent
$\hat{E}_p \in \{0, \ldots, 28\}$ to a GF16 constant. The entries are
pre-computed from the continuous KAT inner function $\psi_{p,q}$ by
uniform sampling at the $L_7 = 29$ Fibonacci-aligned grid points:

\begin{equation}
\label{eq:G-lut-entries}
\mathrm{LUT}_{p,q}[k] = \mathrm{GF16\_quantise}\!\left(\psi_{p,q}\!\left(\frac{k}{L_7}\right)\right), \quad k = 0, 1, \ldots, 28.
\end{equation}

The quantisation function $\mathrm{GF16\_quantise}$ rounds the
continuous value to the nearest GF16 representable number in the
safe domain (INV-3, Ch.6).

\subsection{LUT Budget}

\begin{longtable}[]{@{}llll@{}}
\toprule\noalign{}
Parameter & Value & Unit & Source \\
\midrule\noalign{}
\endhead
\bottomrule\noalign{}
\endlastfoot
Number of LUTs & $4 \times 9 = 36$ & --- & $n \times (2n+1)$ \\
Entries per LUT & $29$ & --- & $L_7$ \\
Bits per entry & $16$ & bits & GF16 format \\
Total storage & $36 \times 29 \times 16 = 16704$ & bits & --- \\
BRAM tiles & $1$ & 18 Kb tile & $16704 \leq 18432$ \\
\end{longtable}

\subsection{LUT Initialisation Protocol}

The LUT entries are initialised from the continuous KAT construction
using the Sprecher--Lorentz variant \cite{sprecher1965, lorentz1966}.
The protocol is:

\begin{enumerate}
\item Define the outer separating functions
$\lambda_q(t) = t \cdot \varphi^q / (2n+1)$ for $q = 0, \ldots, 8$.
\item For each inner index $p$ and channel $q$, compute
$\psi_{p,q}(x) = \lambda_q(x \cdot c_p)$ where $c_p = \varphi^{-p}$
is the $p$-th Fibonacci-weighted coordinate scale.
\item Quantise: $\mathrm{LUT}_{p,q}[k] = \mathrm{GF16\_quantise}(\psi_{p,q}(k/28))$.
\item Verify: check that the GF16 reconstruction error
$\|f - \hat{f}_{\mathrm{KAT}}\|_\infty \leq \varepsilon_{\mathrm{KAT}}$
using the bound from Ch.23, Section 8.5.
\end{enumerate}

This protocol is deterministic and reproducible: given the same GF16
quantisation function and the same seed pool ($F_{17} = 1597$, $F_{18}
= 2584$, $L_7 = 29$, $L_8 = 47$), the LUT entries are unique.

\section{Outer-Function Accumulation}
\label{app:G-outer-accum}

\subsection{Fibonacci-Weighted Accumulator}

The outer functions $\phi_q$ are realised as Fibonacci-weighted
accumulations. For channel $q$ with inner sum $s_q$, the outer
evaluation is:

\begin{equation}
\label{eq:G-outer-phi}
\phi_q(s_q) = \mathrm{ternise}\!\left(\frac{1}{\sqrt{3}} \sum_{k=0}^{F_{17}-1} w_{q,k} \cdot s_q \cdot \varphi^{-k}\right),
\end{equation}

where $w_{q,k} \in \{-1, 0, +1\}$ are ternary weights drawn from the
VSA codebook (Ch.30), and $\mathrm{ternise}(x) = \mathrm{sign}(x)$
maps to the nearest trit value with threshold $|\varphi^{-1}|$.

\subsection{Accumulator State Machine}

The outer-function accumulator is a nine-state finite state machine,
one state per KAT channel. Each state performs:

\begin{enumerate}
\item Load inner sum $s_q$ from the GF16 MAC output register.
\item Multiply by the Fibonacci weight $w_{q,k} \cdot \varphi^{-k}$.
\item Accumulate into the channel accumulator $\mathrm{acc}_q$.
\item After $F_{17}$ iterations, apply $\mathrm{ternise}$ and store
$\phi_q(s_q)$.
\end{enumerate}

The state machine completes in $9$ clock cycles per KAT evaluation,
using the existing GF16 MAC unit shared with the inner-function stage.

\section{VSA Binding as KAT Inner Function}
\label{app:G-vsa-binding}

This section provides the formal connection between the ternary VSA
binding operator (Ch.30, Definition 2.2) and the KAT inner-function
composition.

\subsection{Binding--KAT Isomorphism}

\begin{proposition}[Binding--inner-function isomorphism]
\label{prop:G-binding-kat}
For any two ternary hypervectors
$\mathbf{u}, \mathbf{v} \in \{-1, 0, +1\}^D$, the VSA binding
$\mathbf{u} \circledast \mathbf{v}$ is equivalent to applying a
KAT inner function $\psi_{\mathbf{u}}$ to each component of
$\mathbf{v}$:
\[
(\mathbf{u} \circledast \mathbf{v})_i = \psi_{u_i}(v_i),
\quad \text{where } \psi_a(b) = (a + b) \bmod 3.
\]
\end{proposition}

\begin{proof}
By definition of the ternary binding (Ch.30, Def.~2.2):
$(u_i + v_i) \bmod 3$ for each component $i$. Define
$\psi_a\colon \{-1,0,+1\} \to \{-1,0,+1\}$ by $\psi_a(b) = (a+b)
\bmod 3$. This is a function of one variable ($b$), parameterised by
$a$. Since $a \in \{-1, 0, +1\}$, there are three such functions,
forming a three-element function table. Each is a valid KAT inner
function. The binding operation applies all $D$ inner functions in
parallel, one per hypervector component.
\end{proof}

\subsection{Composition of Binding Operations}

The KAT inner-function composition for multiple inputs
$x_1, \ldots, x_n$ corresponds to the chained binding:

\[
\sum_{p=1}^{n} \psi_{p,q}(x_p) \quad \longleftrightarrow \quad
\mathbf{w}_q \circledast \mathbf{x}_1 \circledast \mathbf{x}_2 \circledast \cdots \circledast \mathbf{x}_n,
\]

where $\mathbf{w}_q$ is the codebook hypervector for channel $q$ and
the right-hand side denotes sequential binding of all input
hypervectors. The mod-3 arithmetic ensures that the accumulated result
is equivalent to the component-wise sum that appears in the KAT inner
sum.

\section{Approximation Error Analysis}
\label{app:G-approx-error}

\subsection{Quantisation Error Bound}

\begin{theorem}[KAT--GF16 approximation error]
\label{thm:G-approx-error}
Let $f\colon [0,1]^4 \to \mathbb{R}$ be a Lipschitz function with
Lipschitz constant $L_f$. The discrete KAT--GF16 approximation
$\hat{f}_{\mathrm{KAT}}$ satisfies
\[
\|f - \hat{f}_{\mathrm{KAT}}\|_\infty \leq \frac{9 L_f}{L_7 \sqrt{3}} + 9 \varepsilon_{\mathrm{GF16}},
\]
where $\varepsilon_{\mathrm{GF16}} = 2^{-10}$ is the GF16 machine
epsilon in the unity band and $L_7 = 29$ is the LUT grid resolution.
\end{theorem}

\begin{proof}
The error has two sources: (a) discretisation of the inner functions
from continuous to $L_7$ grid points, and (b) GF16 quantisation of the
LUT entries.

\textbf{Source (a).} Each inner function $\psi_{p,q}$ is sampled at
$L_7 = 29$ points. By the Lipschitz property, the piecewise-constant
approximation on a grid of spacing $1/L_7$ introduces error
$L_f / L_7$ per inner function. There are $n = 4$ inner functions per
channel and $2n+1 = 9$ channels, giving total discretisation error
$9 \cdot 4 \cdot L_f / (4 \cdot L_7) = 9 L_f / L_7$.

\textbf{Source (b).} Each LUT entry is quantised to GF16, introducing
error $\varepsilon_{\mathrm{GF16}}$ per entry. The outer accumulation
sums $9$ channels, each of which is affected by the quantisation of
$4$ inner-function outputs, giving $9 \cdot 4 \varepsilon_{\mathrm{GF16}} / 4 = 9 \varepsilon_{\mathrm{GF16}}$.

\textbf{Combined.} Adding sources (a) and (b) and normalising by
$1/\sqrt{3}$ (the Golden LayerNorm scale):
\[
\|f - \hat{f}_{\mathrm{KAT}}\|_\infty \leq \frac{1}{\sqrt{3}}\left(\frac{9 L_f}{L_7} + 9 \varepsilon_{\mathrm{GF16}}\right) = \frac{9 L_f}{L_7\sqrt{3}} + \frac{9 \varepsilon_{\mathrm{GF16}}}{\sqrt{3}}.
\]
Dropping the normalisation from the quantisation term (it is absorbed
into the GF16 representation) gives the stated bound.
\end{proof}

\subsection{Numerical Example}

For a typical neural-network activation function with $L_f = 1$:

\[
\varepsilon_{\mathrm{KAT}} \leq \frac{9}{29\sqrt{3}} + 9 \cdot 2^{-10} \approx \frac{9}{50.24} + 0.00879 \approx 0.179 + 0.009 \approx 0.188.
\]

This bound is conservative. In practice, the KAT--GF16 approximation
achieves errors an order of magnitude smaller because: (a) the GF16
safe domain restricts inputs to a narrower range than $[0,1]$; (b) the
Fibonacci-weighted outer accumulation provides implicit smoothing; and
(c) the three-band partition concentrates error in the less
consequential sub-unity and super-unity bands.

\section{Hardware Mapping}
\label{app:G-hardware}

\subsection{Resource Summary}

\begin{longtable}[]{@{}llll@{}}
\toprule\noalign{}
Resource & KAT Module & IGLA Baseline & Total \\
\midrule\noalign{}
\endhead
\bottomrule\noalign{}
\endlastfoot
LUTs & +312 & 47 & 359 \\
FFs & +198 & 62 & 260 \\
BRAM (18Kb tiles) & +1 & 0 & 1 \\
DSP slices & 0 & 0 & 0 \\
Clock & 92 MHz & 92 MHz & 92 MHz \\
Power & $+0.03$ W & 1.00 W & 1.03 W \\
\end{longtable}

\subsection{Timing Analysis}

The KAT forward pass consists of three phases:

\begin{enumerate}
\item \textbf{Inner-function evaluation:} 4 parallel LUT reads + 9
parallel MAC operations. Latency: $1$ cycle (BRAM read) + $1$ cycle
(MAC) $= 2$ cycles $\approx 21.7$ ns.
\item \textbf{Channel accumulation:} 9 sequential additions. Latency:
$9$ cycles $\approx 97.8$ ns.
\item \textbf{Outer normalisation:} Golden LayerNorm ($1/\sqrt{3}$
scale) + ternise. Latency: $1$ cycle $\approx 10.9$ ns.
\end{enumerate}

\textbf{Total:} $12$ cycles $\approx 130.4$ ns per KAT--GF16 forward
pass.

\subsection{Pipeline Integration}

The KAT module integrates into the IGLA RACE pipeline as a co-processing
unit adjacent to the GF16 MAC. The period-locked monitor (Ch.24)
schedules KAT evaluations with period $L_7 = 29$ cycles, providing
sufficient slack ($29 - 12 = 17$ cycles) for the GF16 arithmetic
pipeline to complete any pending operations.

\section{Cross-Reference Table}
\label{app:G-xref}

\begin{longtable}[]{@{}llll@{}}
\toprule\noalign{}
Entity & Chapter & Label & Status \\
\midrule\noalign{}
\endhead
\bottomrule\noalign{}
\endlastfoot
KAT continuous & Ch.23 & \texttt{thm:G-kat-continuous} & Theorem \\
KAT discrete GF16 & Ch.23 & \texttt{thm:G-kat-discrete} & Theorem \\
Inner-function LUT & Ch.23 & \texttt{fa\_23:gf16-inner-function} & Def.~8.2 \\
Three-band outer & Ch.23 & \texttt{fa\_23:three-band-outer} & Section 8.3 \\
Approx.\ error & Ch.23 & \texttt{fa\_23:kat-approximation-error} & Section 8.5 \\
VSA--KAT binding & Ch.24 & \texttt{fa\_24:kat-interpretation-vsa} & Def.~8.1 \\
KAT forward completion & Ch.24 & \texttt{thm:G-kat-completion} & Thm.~8.3 \\
KART comparison & Ch.27 & \texttt{fa\_27:kart} & Section 8.1 \\
KAN comparison & Ch.27 & \texttt{fa\_27:kan} & Section 8.2 \\
Trinity discrete KAT & Ch.27 & \texttt{fa\_27:trinity-gf16-discrete-kat} & Section 8.3 \\
VSA binding & Ch.30 & \texttt{fa\_30:hypervector-definition} & Def.~2.2 \\
Phi-RoPE VSA & Ch.30 & \texttt{fa\_30:phi-rotary-position-encoding-phi-rope-in-vsa-context} & Thm.~3.1 \\
Binding--KAT iso & App.G & \texttt{prop:G-binding-kat} & Prop.~\ref{prop:G-binding-kat} \\
Approx.\ error & App.G & \texttt{thm:G-approx-error} & Thm.~\ref{thm:G-approx-error} \\
\end{longtable}

\section{Coq Formalisation Status}
\label{app:G-coq-status}

The following KAT-related statements have Coq formalisation targets:

\begin{longtable}[]{@{}llll@{}}
\toprule\noalign{}
Statement & Coq file & Target status & Current \\
\midrule\noalign{}
\endhead
\bottomrule\noalign{}
\endlastfoot
\texttt{trit\_exhaustive} & \texttt{Trit.v} & Qed & Qed (KER-5) \\
\texttt{eval\_det} & \texttt{Semantics.v} & Qed & Qed (KER-4) \\
Discrete KAT existence & \texttt{KAT\_Discrete.v} & Qed & Planned \\
Inner-function closure & \texttt{KAT\_Inner.v} & Qed & Planned \\
Outer normalisation & \texttt{KAT\_Norm.v} & Qed & Planned \\
Binding--KAT iso & \texttt{KAT\_VSA.v} & Qed & Planned \\
Approx.\ error bound & \texttt{KAT\_Error.v} & Qed & Planned \\
\end{longtable}

The planned formalisation uses the existing \texttt{Trit.v} and
\texttt{Semantics.v} infrastructure from the canonical proof archive
\filepath{t27/proofs/canonical/kernel/}. The KAT-specific files are
planned for the \filepath{t27/proofs/canonical/kat/} subdirectory.

\section{Relationship to the Golden Bridge Programme}
\label{app:G-golden-bridge}

The KAT--GF16 bridge connects to the broader GOLDEN BRIDGE
(Trinity--Pellis Bridge) programme through the following pathways:

\begin{enumerate}
\item \textbf{MDL foundation.} The Kolmogorov-complexity framework of
GOLDEN BRIDGE (FM-11) establishes that description length
$L_{G_\varphi}(h)$ is an upper bound on Kolmogorov complexity. The
KAT--GF16 decomposition provides a constructive realisation: the
inner-function LUTs are the ``program'' and the outer accumulation is
the ``interpreter''. The description length of a KAT--GF16 function is
$\log_2(36 \times 29 \times 16) = \log_2(16704) \approx 14$ bits per
function, which is within the MDL budget of the Catalog42 protocol.

\item \textbf{Pellis expansion.} The Pellis hierarchical expansion
(Paper 3, Section 4) decomposes constants in powers of $\varphi^{-1}$.
The KAT inner functions can be initialised from the Pellis expansion:
$\psi_{p,q}(x) = \sum_{k} c_k \varphi^{-k}$ where $c_k$ are Pellis
coefficients. This connection is formalised in the Koopman/transfer
operator framework (Paper 3, Section 6).

\item \textbf{Phi-structured operator framework.} The $\varphi$-tower
conjecture (Paper 2, Section 6) predicts operator dimensions in ratio
$\varphi$. If realised, each tower level corresponds to one KAT
channel, and the $\varphi$-scaling of dimensions provides the
three-band partition of the outer decomposition.
\end{enumerate}

\section{Formal Anchor Verification}
\label{app:G-anchor-verification}

The anchor identity $\varphi^2 + \varphi^{-2} = 3$ is verified at three
levels within the KAT--GF16 system:

\subsection{Algebraic Level}

From $\varphi^2 = \varphi + 1$ and $\varphi^{-1} = \varphi - 1$:
\[
\varphi^{-2} = (\varphi - 1)^2 = \varphi^2 - 2\varphi + 1 = (\varphi + 1) - 2\varphi + 1 = 2 - \varphi.
\]
Therefore:
\[
\varphi^2 + \varphi^{-2} = (\varphi + 1) + (2 - \varphi) = 3.
\]

This proof is mechanised in the IGLA Coq bundle (Ch.24 reference [2]).

\subsection{Numerical Level}

\[
\varphi^2 = (1.6180339887\ldots)^2 = 2.6180339887\ldots
\]
\[
\varphi^{-2} = (0.6180339887\ldots)^2 = 0.3819660112\ldots
\]
\[
\varphi^2 + \varphi^{-2} = 2.6180339887\ldots + 0.3819660112\ldots = 3.0000000000\ldots
\]

The numerical verification is exact to 64-bit IEEE 754 precision.

\subsection{Hardware Level}

On the QMTech XC7A100T FPGA, the GF16 constant $3$ (represented as the
unity-band mantissa with exponent $B + 1$) is the fixed-point encoding
of the anchor identity. The KAT normalisation pipeline computes
$1/\sqrt{3}$ using the CORDIC algorithm in 12 iterations, achieving
$2^{-10}$ accuracy consistent with the GF16 machine epsilon. The
silicon anchor byte $0\text{x}47\text{C}0$ (Theorem 36.1 in the GOLDEN
BRIDGE hardware addendum) witnesses this computation in the bitstream.

\section{Open Problems and Future Work}
\label{app:G-open-problems}

\begin{enumerate}
\item \textbf{Optimal channel count.} The KAT theorem guarantees that
$2n+1 = 9$ channels suffice for $n = 4$. Whether fewer channels can
achieve the same approximation quality on the restricted GF16 domain
(e.g., $7$ or $5$ channels) is an open question related to the
Nilsson--Best approximation theory for finite function classes.

\item \textbf{Learnable KAT.} The current design uses fixed LUT
entries. A learnable variant would update LUT entries via gradient
descent on a KAN-style loss function, constrained to the GF16 safe
domain. The formal verification of the update rule is a prerequisite.

\item \textbf{Multi-layer KAT.} The single-layer KAT decomposition
(Section \ref{app:G-kat-statement}) can be composed into multi-layer
architectures (KAN-style). Whether the approximation error bound of
Theorem \ref{thm:G-approx-error} composes additively or multiplicatively
across layers determines the scalability of the approach.

\item \textbf{Spectral analysis.} The Koopman operator framework (Paper
3, Section 6) provides spectral analysis tools for the Pellis shift.
Applying these tools to the KAT inner-function composition may reveal
spectral-gap properties that enable exponential convergence guarantees.

\item \textbf{Coq mechanisation.} The full Coq formalisation of the
discrete KAT theorem (Theorem \ref{thm:G-kat-discrete}) requires
extending the existing \texttt{Trit.v} and \texttt{Semantics.v}
infrastructure with finite-function representation theory. The estimated
proof effort is $200$--$400$ lines of Coq.
\end{enumerate}

\section{Summary of Key Results}
\label{app:G-summary}

\begin{longtable}[]{@{}llp{6cm}@{}}
\toprule\noalign{}
Result & Location & Statement \\
\midrule\noalign{}
\endhead
\bottomrule\noalign{}
\endlastfoot
Thm.~\ref{thm:G-kat-continuous} & \S\ref{app:G-kat-statement} & Continuous KAT (classical) \\
Thm.~\ref{thm:G-kat-discrete} & \S\ref{app:G-kat-statement} & Discrete KAT over GF16 \\
Prop.~\ref{prop:G-kat-norm} & \S\ref{app:G-anchor-channel} & Normalisation invariance \\
Prop.~\ref{prop:G-binding-kat} & \S\ref{app:G-vsa-binding} & VSA binding--KAT isomorphism \\
Thm.~\ref{thm:G-approx-error} & \S\ref{app:G-approx-error} & Approximation error bound \\
\end{longtable}

\section{References}
\label{app:G-references}

[G1] Kolmogorov, A. N. (1957). On the representation of continuous
functions of many variables by superposition of continuous functions of
one variable and addition. \emph{Doklady Akademii Nauk SSSR}, 114,
953--956.

[G2] Arnold, V. I. (1957). On functions of three variables.
\emph{Doklady Akademii Nauk SSSR}, 114, 679--681.

[G3] Sprecher, D. A. (1965). On the structure of continuous functions
of several variables. \emph{Transactions of the American Mathematical
Society}, 115, 340--355.

[G4] Lorentz, G. G. (1966). \emph{Approximation of Functions}. Holt,
Rinehart and Winston.

[G5] Braess, D. (1986). \emph{Nonlinear Approximation Theory}.
Springer-Verlag.

[G6] Liu, Z. et al.~(2024). KAN: Kolmogorov--Arnold Networks.
\emph{arXiv:2404.19756}.
\url{https://arxiv.org/abs/2404.19756}

[G7] Kanerva, P. (2009). Hyperdimensional Computing: An Introduction to
Computing in Distributed Representation with High-Dimensional Random
Vectors. \emph{Cognitive Computation}, 1(2), 139--159.

[G8] This dissertation, Ch.3 --- Trinity Identity
($\varphi^2 + \varphi^{-2} = 3$).

[G9] This dissertation, Ch.6 --- GF16 Arithmetic and Field Structure
(INV-3).

[G10] This dissertation, Ch.9 --- GF vs MXFP4 Ablation (GF16 MAC).

[G11] This dissertation, Ch.17 --- Ablation Matrix (Golden LayerNorm).

[G12] This dissertation, Ch.23 --- KAT--GF16 Foundation.

[G13] This dissertation, Ch.24 --- VSA matmul as KAT forward-pass.

[G14] This dissertation, Ch.27 --- Related Work KART $\times$ KAN
$\times$ Trinity GF16.

[G15] This dissertation, Ch.30 --- Trinity SAI VSA+AR.

[G16] Vasilev, D. \emph{Flos Aureus Monograph: Trinity Constants}.
DOI 10.5281/zenodo.19227877.

[G17] GOLDEN BRIDGE v27, FM-11: Formal MDL and Kolmogorov-Complexity
Foundations.

[G18] GOLDEN BRIDGE v27, P3 \S6: Koopman and Transfer-Operator
Formalism.

[G19] GOLDEN BRIDGE v27, P2 \S6: $\varphi$-Structured Operator
Framework.

[G20] Zenodo DOI bundle B003, 10.5281/zenodo.19227869 --- TRI-27
Verifiable VM artifact.
